\chapter*{术语表}
\addcontentsline{toc}{chapter}{术语表}

{\small
\begin{longtable}{
  >{\raggedright\arraybackslash}p{0.16\linewidth}
  >{\raggedright\arraybackslash}p{0.28\linewidth}
  >{\raggedright\arraybackslash}p{0.42\linewidth}
}
\toprule
中文 & 英文 & 简要解释 \\
\midrule
翻译单元 & translation unit & 一个源文件经过预处理后的编译输入。 \\
中间表示 & intermediate representation & 编译器用于优化和 lowering 的内部表示。 \\
抽象语法树 & abstract syntax tree, AST & 表示源码语法结构的树，例如声明、语句和表达式的层级关系。 \\
控制流图 & control-flow graph, CFG & 把基本块和跳转边连起来的图，用于分析分支、循环和优化。 \\
LLVM & compiler infrastructure & 编译器基础设施项目，本书常用 LLVM IR 和 Clang 诊断观察优化。 \\
GCC & GNU Compiler Collection & GNU 编译器集合，本书常用 GIMPLE 或优化诊断与 LLVM 对照。 \\
Clang & LLVM C/C++ frontend & LLVM 生态中的 C/C++ 前端，用于生成 IR、汇编、诊断和编译数据库。 \\
SSA & static single assignment & 静态单赋值形式，每个虚拟值只定义一次，便于数据流分析和优化。 \\
GIMPLE & GCC intermediate representation & GCC 使用的一类中间表示，常用于观察优化前后的树形/SSA 事实。 \\
LTO & link-time optimization & 链接期优化，允许编译器跨翻译单元做内联、去虚化和全局分析。 \\
DCE & dead code elimination & 死代码消除，删除无使用且无可观察副作用的计算。 \\
LICM & loop-invariant code motion & 循环不变代码外提，把安全的不变计算移到循环外。 \\
mem2reg & memory to register promotion & 把没有地址逃逸的栈槽提升成 SSA 值，是理解优化 IR 的入口。 \\
ASM & assembly & 汇编文本或汇编语言层表示，用助记符表达机器指令。 \\
GAS & GNU assembler & GNU 汇编器和相关语法生态，Linux 下常见。 \\
MASM & Microsoft Macro Assembler & Windows/MSVC 生态常见汇编器。 \\
MSVC & Microsoft Visual C++ & 微软 C++ 编译器和工具链。 \\
COFF & Common Object File Format & 一类目标文件格式族，Windows PE/COFF 生态中尤其常见。 \\
SEH & structured exception handling & Windows 结构化异常处理机制。 \\
FFI & foreign function interface & 跨语言调用边界，结构体布局、调用约定和异常传播都要明确。 \\
应用二进制接口 & application binary interface, ABI & 函数调用、寄存器、栈、对象布局和符号链接的二进制合同。 \\
应用程序接口 & application programming interface, API & 源码层调用合同；API 能编译不代表 ABI 一定兼容。 \\
指令集架构 & instruction set architecture, ISA & 软件可见的指令、寄存器和状态集合。 \\
x86-64 & 64-bit x86 architecture & 64 位 x86 架构名称，常和 AMD64 指同一类用户态目标。 \\
AMD64 & 64-bit x86 architecture & AMD 最早推动的 64 位 x86 扩展名，System V AMD64 ABI 采用这个命名。 \\
IA-32 & 32-bit Intel architecture & 32 位 x86 架构名，和 64 位 x86-64/AMD64 不是同一 ABI 边界。 \\
RISC-V & open instruction set architecture & 开放指令集架构。它和 x86-64 的寄存器、指令编码、ABI 和内存模型都不同。 \\
操作系统 & operating system, OS & 管理进程、线程、内存、文件和设备的系统软件。 \\
PID & process id & 进程编号，用来在 \code{/proc}、调试器、日志和系统工具中定位一个运行实例。 \\
PC & program counter & 程序计数器，保存下一条将要取指的指令地址。x86-64 文本中常通过 RIP 观察它。 \\
VM & virtual memory & 虚拟内存，进程看到的地址空间由操作系统和页表映射到物理内存、文件或设备。 \\
PROC & proc filesystem & Linux 的 \code{/proc} 伪文件系统，把进程、内核和硬件状态暴露成可读文件。 \\
PR & pull request & 代码合并请求。底层或性能 PR 应附带二进制、反汇编、benchmark 和环境证据。 \\
ID & identifier / identity & 身份或编号，用于唯一定位用户、进程、实验样本或资源对象。 \\
UNIX & operating system family and tradition & 一类操作系统传统和接口生态，Linux 继承了许多 UNIX/POSIX 的进程、文件、信号和管道思想。 \\
EOF & end of file & 输入结束或文件读完状态。它不是普通字符，而是读取结果或流状态。 \\
RLIMIT & resource limit & 进程资源限制，例如打开文件数、栈大小或地址空间上限。 \\
RSS & resident set size & 常驻内存大小，表示当前实际驻留在物理内存中的页面量。 \\
CI & continuous integration & 持续集成，用自动构建、测试和检查保护每次提交。 \\
I/O & input/output & 输入输出，包含文件、终端、网络和设备访问；计时实验必须说明是否包含它。 \\
STL & Standard Template Library & C++ 标准库容器、迭代器和算法的常见统称。 \\
CC & C compiler variable & 常见构建环境变量，用于指定 C 编译器。 \\
CXX & C++ compiler variable & 常见构建环境变量，用于指定 C++ 编译器。 \\
CXXFLAGS & C++ compiler flags & 常见构建环境变量，用于指定 C++ 编译选项。 \\
C++20 & C++ standard version & C++ 标准版本名。第一册示例使用它时，构建命令必须显式启用对应标准。 \\
O0/O2/O3 & optimization levels & 编译优化级别。O0 便于调试，O2/O3 更接近性能构建，但会改变代码形态和调试体验。 \\
DNDEBUG & disable debug assertions & 常用于关闭 \code{assert} 的预处理宏，发布构建中会改变断言相关行为。 \\
ELF & Executable and Linkable Format & Linux 常见目标文件、共享库和可执行文件格式。 \\
ELF64 & 64-bit ELF & 64 位 ELF 文件类别。 \\
INTERP & program interpreter entry & ELF 中指定动态加载器路径的字段。 \\
NEEDED & required shared library entry & ELF 动态段中记录需要加载哪些共享库的字段。 \\
PLT/GOT & Procedure Linkage Table / Global Offset Table & 动态链接中定位外部函数和全局对象的表结构。 \\
GOTPCREL & GOT PC-relative relocation & x86-64 位置无关代码中相对 RIP 访问 GOT 项的常见重定位或寻址记号。 \\
JUMP\_SLOT & dynamic call relocation & ELF 动态链接中用于函数调用目标绑定的重定位类型之一。 \\
DWARF & Debugging With Attributed Record Formats & 常见调试信息格式，记录源码行号、类型、变量位置和栈展开信息。 \\
GDB & GNU Debugger & GNU 调试器，用 DWARF 等信息支持断点、单步、查看变量和调用栈。 \\
ASan & AddressSanitizer & 地址消毒器，用于发现越界、use-after-free 等内存错误。 \\
UBSan & UndefinedBehaviorSanitizer & 未定义行为消毒器，用于发现有符号溢出、非法移位、错误对齐等问题。 \\
TSan & ThreadSanitizer & 线程消毒器，用于发现数据竞争和部分同步错误。 \\
CFI & call frame information & 调用帧信息，帮助调试器、异常处理和采样器从当前栈帧回到调用者。 \\
RPATH & runtime library search path & ELF 动态段中的库搜索路径字段，会影响运行时实际加载哪个共享库。 \\
RUNPATH & runtime library search path & ELF 动态段中的另一类运行时库搜索路径字段，现代链接器更常见。 \\
LD\_LIBRARY\_\allowbreak PATH & dynamic linker search path & 运行时动态库搜索路径环境变量，可能改变程序行为和崩溃位置。 \\
LD\_DEBUG & dynamic linker debug output & 动态链接器调试环境变量，可输出库查找、符号绑定和重定位过程，常用于排查加载到哪个共享库。 \\
RIP/RFLAGS & instruction pointer / flags register & RIP 指向下一条指令；RFLAGS 保存条件码、进位、方向等状态位。 \\
EFLAGS & x86 flags register & 32 位 x86 标志寄存器名；x86-64 中 RFLAGS 是其 64 位扩展。 \\
ZF/CF/SF/PF & zero/carry/sign/parity flags & ZF 表示零；CF 表示进位或借位；SF 表示符号位；PF 表示低字节奇偶。 \\
OF & overflow flag & 有符号溢出标志，常和 SF、ZF 一起用于有符号比较后的条件跳转。 \\
SIB & scale-index-base & x86 寻址编码字段，用 base、index、scale 组合出有效地址。 \\
LEA & load effective address & x86 地址计算指令，常被编译器用作不访存的整数加法和乘加组合。 \\
REX & x86-64 instruction prefix & x86-64 扩展前缀，用来选择扩展寄存器或 64 位操作数。 \\
VEX/EVEX & vector instruction prefixes & x86 向量扩展前缀，用来表达更宽向量、更多寄存器和更丰富操作数。 \\
XMM/YMM & vector registers & x86 SIMD 寄存器族。XMM 是 128 位，YMM 是 256 位。 \\
CPL & current privilege level & 当前特权级，区分用户态、内核态等执行权限。 \\
DPL & descriptor privilege level & 描述符特权级，参与 x86 权限检查。 \\
RPL & requested privilege level & selector 低位携带的请求特权级，参与部分 x86 权限检查。 \\
IDT & interrupt descriptor table & 中断描述符表，处理器用它定位中断和异常处理入口。 \\
TSS & task state segment & x86 中保存特权级切栈等信息的结构，现代系统仍用于异常和中断入口。 \\
RSP0 & ring-0 stack pointer & TSS 中常见的内核栈指针字段，用于从用户态进入 ring 0 时切栈。 \\
IST & interrupt stack table & 中断栈表，为危险异常提供专用栈入口。 \\
MSR & model-specific register & 模型专用寄存器，系统调用入口、控制位和硬件特性常通过它配置。 \\
IA32 & Intel architecture 32 prefix & x86 体系中沿用的寄存器命名前缀，例如 \code{IA32\_LSTAR}。 \\
IPI & inter-processor interrupt & 处理器间中断，一个核心向另一个核心发送的中断事件。 \\
\#PF & page fault exception & 页故障异常，表示地址翻译、页存在性或权限检查失败。 \\
CET & control-flow enforcement technology & 控制流保护技术，例如 shadow stack，用于降低返回地址和间接跳转被篡改的风险。 \\
XD & execute disable & 不可执行页权限的另一种常见名称，和 NX 表达同类取指禁止语义。 \\
SIGSEGV & segmentation violation signal & 非法内存访问信号，常见于空指针、越界、权限错误或损坏地址。 \\
SIGPIPE & broken pipe signal & 向已关闭的管道或 socket 写入时常见的信号。 \\
SIGTERM & termination signal & 请求进程终止的常见信号，程序可捕获后做清理。 \\
SIGINT & interrupt signal & 终端中断信号，常由 Ctrl-C 触发。 \\
SIGABRT & abort signal & 进程主动异常终止信号，常见于 \code{abort} 或断言失败。 \\
SIGALRM & alarm signal & 定时器告警信号，会打断阻塞或普通控制流。 \\
SIGSTKSZ & signal stack size constant & POSIX 信号备用栈相关常量，用于 \code{sigaltstack} 这类场景。 \\
SA\_RESTART & restart interrupted syscall flag & POSIX 信号处理标志，允许部分被信号打断的系统调用自动重启。 \\
EINTR & interrupted system call & 系统调用被信号中断。可靠系统代码要明确是否重试。 \\
ENOENT & no such file or directory & 常见 errno，表示路径不存在。 \\
EAGAIN & try again & 常见 errno，表示资源暂时不可用或非阻塞操作需要稍后重试。 \\
EMFILE & too many open files & 进程打开文件描述符达到资源限制。 \\
POSIX & portable operating system interface & Unix-like 系统接口规范，信号、文件描述符和线程接口常以它为边界。 \\
VFS & virtual file system & Linux 虚拟文件系统层，把不同文件系统统一到文件操作接口。 \\
PTE & page table entry & 页表项，记录虚拟页到物理页的映射和权限位。 \\
线程局部存储 & thread-local storage, TLS & 每个线程拥有独立变量实例的存储机制。 \\
地址空间布局随机化 & address space layout randomization, ASLR & 运行时随机化关键地址，降低地址预测攻击风险。 \\
位置无关代码 & position-independent code, PIC/PIE & 能在不同装载地址运行的代码或可执行文件。 \\
RTTI & run-time type information & C++ 运行期类型信息，支持 \code{dynamic_cast} 和 \code{typeid} 等机制。 \\
LSDA & language-specific data area & 异常展开表中的语言特定数据区域，用来描述 catch、cleanup 和 landing pad。 \\
架构状态 & architectural state & ISA 规定的软件可见状态，如寄存器、内存、flags。 \\
微架构 & microarchitecture & CPU 内部实现 ISA 的方式，如 pipeline、cache、uops。 \\
ALU & arithmetic logic unit & 算术逻辑单元，执行整数算术、逻辑和比较等操作。 \\
ROB & reorder buffer & 重排序缓冲区，保存乱序执行但尚未按程序顺序提交的微操作。 \\
BTB & branch target buffer & 分支目标缓存，用来预测跳转目标地址。 \\
RSB & return stack buffer & 返回栈缓冲，用来预测函数返回地址。 \\
RISC & reduced instruction set computer & 精简指令集计算机风格的架构类别，通常强调较规则的指令编码和 load/store 结构。 \\
SIMD & single instruction, multiple data & 一条指令处理多个数据 lane 的执行方式。 \\
SSE & Streaming SIMD Extensions & x86 较早的 SIMD 扩展族，常作为理解向量寄存器和 lane 的入口。 \\
AVX/AVX2 & Advanced Vector Extensions & x86 向量指令扩展，用于提升数据并行吞吐。 \\
AVX-512 & Advanced Vector Extensions 512-bit & x86 512 位向量扩展，提供更宽寄存器和 mask 机制，但需要验证频率和平台覆盖。 \\
ZMM & 512-bit vector register & AVX-512 使用的 512 位向量寄存器名。 \\
FMA & fused multiply-add & 一条指令完成乘加，并减少中间舍入。 \\
TLB & translation lookaside buffer & 地址翻译缓存，保存虚拟页到物理页映射。 \\
L1/L1D/L2/L3 & cache hierarchy & CPU 缓存层级。L1D 是一级数据缓存；L2 居中；L3 通常较大并可能被多个核心共享。 \\
延迟 & latency & 操作输入 ready 到输出 ready 的时间。 \\
吞吐 & throughput & 稳态下单位时间能完成的操作数量。 \\
PMU & performance monitoring unit & 提供硬件性能计数器的单元，用于观察周期、cache miss、分支失败等事件。 \\
TSC & time stamp counter & 处理器时间戳计数器，常用于低开销计时，但要理解频率、迁核和序列化边界。 \\
IPC & instructions per cycle & 每周期执行指令数，是观察 CPU 利用情况的指标之一。 \\
ILP & instruction-level parallelism & 指令级并行，互不依赖指令可以重叠执行的能力。 \\
SRAM & static random-access memory & 静态随机访问存储器，速度快、面积成本高，常用于 CPU cache。 \\
ULP & unit in the last place & 浮点末位单位，用来描述相邻可表示浮点数之间的误差尺度。 \\
CSV & comma-separated values & 逗号分隔文本表格，常用于保存 benchmark 和 PMU 采样结果。 \\
GB/s & gigabytes per second & 每秒十亿字节量级的数据吞吐，报告时要说明读、写还是读写合计。 \\
GFLOP/s & giga floating-point operations per second & 每秒十亿次浮点运算，必须绑定 FLOP 计数公式。 \\
QPS & queries per second & 每秒请求数，服务端吞吐指标；必须和延迟分位、错误率和资源占用一起解释。 \\
LLC & last-level cache & 最后一级缓存，通常是多个核心共享的大缓存。 \\
NUMA & non-uniform memory access & 非统一内存访问，不同内存节点访问成本不同。 \\
SMT & simultaneous multithreading & 同一个物理核心同时运行多个硬件线程，共享部分执行和缓存资源。 \\
COW & copy-on-write & 写时复制，多个对象先共享底层页或缓冲，写入时再复制出私有版本。 \\
RMW & read-modify-write & 读改写操作，把读取旧值、计算新值、写回合成一个不可分割的原子更新序列。 \\
JIT & just-in-time compilation & 运行时即时编译或生成代码，常见于虚拟机、正则、解释器和插件系统。 \\
UAF & use-after-free & 释放后继续使用对象或内存，是常见内存安全错误。 \\
ROP & return-oriented programming & 利用被破坏的返回地址拼接已有指令片段的攻击技术，是理解 CFI、CET 和 shadow stack 的背景。 \\
RTL & register transfer language & GCC 后端常见中间表示层，接近寄存器传送和机器指令选择。 \\
HPC & high-performance computing & 高性能计算，关注吞吐、带宽、并行性和可复现实验。 \\
GEMM & general matrix-matrix multiplication & 矩阵乘矩阵，高算术强度计算内核的代表。 \\
VNNI & vector neural network instructions & 面向低精度 dot-product 的向量指令能力。 \\
FTZ/DAZ & flush-to-zero / denormals-are-zero & 浮点非正规数处理模式，会影响性能和数值语义。 \\
MXCSR & SSE control/status register & x86 SSE/AVX 浮点控制状态寄存器，包含舍入、异常和 FTZ/DAZ 等位。 \\
\makecell[l]{ABI 分类\\类别} & ABI classification classes & System V AMD64 ABI 中聚合参数/返回值分类用的类别名，例如 INTEGER、SSEUP、MEMORY、X87；它们决定对象放进通用寄存器、向量寄存器还是内存。 \\
NX & no-execute & 不可执行页权限，阻止把普通数据页当作代码执行。 \\
BSS & block started by symbol & 未初始化全局/静态数据段，装载后通常清零。 \\
VMA & virtual memory area & 进程地址空间中的连续虚拟内存区域。 \\
CFA & canonical frame address & 调用帧地址，DWARF 栈展开中用于定位调用者帧。 \\
ODR & one definition rule & C++ 单一定义规则，要求同一实体在程序中满足一致定义。 \\
IDE & integrated development environment & 集成开发环境，通常包含编辑、构建、调试和代码导航。 \\
README & project readme & 项目说明文件，用于记录构建、运行、依赖和维护入口。 \\
JSON & JavaScript Object Notation & 结构化文本格式，常用于编译数据库、配置和实验报告。 \\
POSIX/ELF 宏 & system/interface macros & 例如 RLIMIT\_NOFILE、AT\_FDCWD、WIFEXITED、WEXITSTATUS；它们是系统接口约定的一部分，不是普通业务变量。 \\
benchmark 常量 & benchmark parameters & 例如 BENCHMARK\_WARMUP\_REPETITIONS、BENCHMARK\_INPUT\_SIZE、P95\_NUMERATOR，用来固定实验口径。 \\
汇编局部标签 & assembler labels & 例如 LJTI0\_0、LBB0\_1、FUNC、ENDP，用于跳转表、基本块或平台汇编标记。 \\
LCG & linear congruential generator & 线性同余伪随机生成器，常用于可复现但不追求高随机质量的 workload。 \\
BLAS & Basic Linear Algebra Subprograms & 基础线性代数接口族，很多高性能库以它为通用入口。 \\
KV cache & key/value cache & Transformer 解码保存历史 K/V 的缓存，本册作为后续 AI runtime 例子出现。 \\
TTFT & time to first token & 首 token 延迟，模型服务冷启动、预热、prefill 和 decode 分层时会用到。 \\
WSL2 & Windows Subsystem for Linux 2 & Windows 上的 Linux 虚拟化环境，性能和 PMU 可用性需单独说明。 \\
UI & user interface & 用户界面。profiling UI 或 IDE 能提高效率，但不能替代可复查命令和原始证据。 \\
RCU & read-copy-update & Linux 内核常见同步思想，读路径尽量无锁，更新路径复制并延迟回收旧版本。 \\
IEEE 754 & floating-point standard & 常见浮点数标准，规定舍入、NaN、Inf、signed zero 等语义。 \\
算术强度 & arithmetic intensity & FLOPs 与数据搬运字节数的比值。 \\
RAII & resource acquisition is initialization & 用对象生命周期管理资源的 C++ 惯用法。 \\
UB & undefined behavior & 未定义行为，触发后标准不再规定程序结果。 \\
DMA & direct memory access & 设备直接访问内存而不需要 CPU 逐字节搬运。 \\
DRAM & dynamic random-access memory & 主存常用的动态随机访问存储器。 \\
\bottomrule
\end{longtable}
}
